\chapter{具身接口层：观测绑定、现实校准与闭环控制}
\label{chp:microscopic_layer_ch}

微观层是智能系统与环境直接交换观测、动作和资源的具身接口。本章不再把它描述为物理世界与某种特殊认知场之间的全息切面，也不由狄拉克源、量子坍缩或热力学隐喻证明感知机制。我们把它重建为一组具有明确数据类型、时间尺度、误差模型和安全边界的闭环算子：上行链路从传感器数据中估计对象、属性、关系与不确定性，下行链路把任务目标变成可执行的参考状态、控制参数和动作约束，反馈链路再用现实结果修订当前状态与长期模型。局部结构网络负责保持对象身份、空间关系和事件来源，全局分布表示负责表达候选解释、置信与连续控制状态；二者只有经过绑定、耦合和读出接口才能共同形成可行动的认知状态。微观层因此不是被动的输入输出端口，而是现实校准、复杂度限流和动作资格审查的第一道闭环。本章保留VTE、源流、形质双通道、现实锚定、共振/投影、效应器和阻抗控制等历史问题，但将其分别落实为可测试的编码、估计、门控、采样和控制模型。

\section{观测编码接口：从原始信号到同源绑定记录}
\label{sec:vte_as_wave_source}

我们把历史术语“变分拓扑编码器”保留为观测编码接口 VTE，但重新规定其对象。设外部环境状态为 $x_t^{e}\in\mathcal{X}_e$，传感器在外部时钟 $t$ 上输出 $y_t\in\mathbb{R}^{m}$；$y_t$ 的各维带有可校准单位，例如像素亮度、帕斯卡、弧度每秒或牛顿。观测模型写为
\[
y_t=h_{\phi_t}(x_t^{e},c_t)+\nu_t,
\]
其中 $c_t$ 是曝光、姿态、温度等传感器条件，$\nu_t$ 是在声明分布或有界集合中的噪声，$\phi_t$ 是可随标定更新的参数。接口的任务不是把“物理能量”变成“语义能量”，而是在有限分辨率和延迟下构造记录 $o_t=(F_t,Q_t,B_t,S_t)$：$F_t$ 是对象位置、姿态、边界和关系候选等结构特征，$Q_t$ 是颜色、纹理、声音、压力等属性特征，$B_t$ 是结构与属性的同源绑定，$S_t$ 是来源、时间戳、标定版本和不确定性。

形与质双通道因此是数据分工，不是玻色子与费米子的物理划分。结构编码器 $E_f:y_{t-w:t}\mapsto F_t$ 从时间窗中估计可达位置、对象轮廓、光流、接触法线和拓扑关系；属性编码器 $E_q:y_{t-w:t}\mapsto Q_t$ 估计颜色分布、材料类别、声音频谱或触觉纹理。二者不能彼此独立后任意组合，而要通过共享感受野、对象身份和采集时间建立绑定。设结构候选为 $f_i\in\mathbb{R}^{d_f}$，属性候选为 $q_j\in\mathbb{R}^{d_q}$，来源兼容度为 $a_{ij}$，则合法绑定集合为
\[
\mathcal{B}_t=\{(i,j):a_{ij}\geq\alpha,\;\Delta t_{ij}\leq\delta_t,\;\text{source}(i)\sim\text{source}(j)\}.
\]
在该集合上才计算绑定概率。红点位于 $(10,20)$、白点位于 $(30,50)$ 的例子中，位置与颜色都携带同一对象索引和时间窗；若红色来自后一对象，系统必须把错配保留为竞争解释，而不是用局部张量积宣称绝无泄漏。

VTE 的“变分”含义是求解带约束估计，而不是服从最小作用量。令解码器 $D$ 将候选记录重构为传感器预测，估计目标可以写为
\[
L_t=\|W_t^{1/2}(y_t-D(F_t,Q_t,B_t,c_t))\|_2^2
+\lambda_f R_f(F_t)+\lambda_qR_q(Q_t)+\lambda_bR_b(B_t),
\]
其中 $W_t$ 依据传感器协方差设定，$R_f$ 约束结构的时间连续性，$R_q$ 抑制不可靠属性，$R_b$ 惩罚一对多错配。所有项先声明为无量纲评分；真实焦耳能耗若纳入优化，需经独立测量和换算。损失下降只能说明当前模型更好重构输入，不说明识别对象真实类别。遮挡、镜面反射、传感器饱和和域外信号都可能产生多个近似最优解，因此接口输出必须包含候选分布而不是单一“本体论波包”。

局部结构网络接收经验证的对象节点、关系边、来源与版本，全局分布状态接收连续特征和候选权重。绑定接口把两者联合为带身份的观测包，读出接口再针对任务选择分辨率。机器人抓取杯子时需要位置、姿态、杯柄关系和摩擦估计；颜色只在区分目标时进入读出。听觉定位则可能让时间差承担结构作用、频谱承担属性作用。我们不从人脑背侧/腹侧通路直接推出这套软件结构，只把神经科学现象视为“不同特征与任务接口可分工”的经验启发。VTE是否有效，最终由身份错配率、定位误差、属性校准、重构残差、处理延迟和下游任务表现共同检验。

接口还必须处理“没有形成对象”的输入。雾、眩光、持续噪声或一片无法分割的纹理可能只支持区域级统计，而不支持稳定节点；此时VTE应输出未绑定特征、失败原因和有效视域，不能为了满足下游模式而虚构对象身份。相反，稀疏事件也可能在时间上累积为对象，例如只有数次触觉接触时，系统可维护暂定身份并等待主动探测。我们据此把编码成功分成特征可测、候选可分、身份可持和任务可读四级，分别报告失败，而不以单一准确率掩盖接口究竟在哪一级失效。

\section{现实校准：观测似然、边界置信与失锚诊断}
\label{sec:section_5326}

现实锚定不是把内部状态硬锁定为传感器数值。传感器可能坏、被欺骗、延迟或只看到环境的一小部分，内部预测也可能提供遮挡补全和噪声抑制。我们把锚定定义为在来源可追踪的观测模型下更新环境信念。令先验为 $p_{t|t-1}(x)$，观测似然为 $p(y_t\mid x,c_t,s_t)$，其中 $s_t$ 记录传感器身份和健康状态，则
\[
p_{t|t}(x)=\frac{p(y_t\mid x,c_t,s_t)p_{t|t-1}(x)}{\int p(y_t\mid x',c_t,s_t)p_{t|t-1}(x')\,dx'}
\]
仅在分母存在且非零时成立。高精度观测使后验更多受似然约束，低精度或异常观测应降低权重，而不是令耦合刚度趋于无穷。相机看到红色并不保证对象“本来就是红色”；曝光、光源和材料模型共同决定似然。现实校准的原则是证据有资格改变信念，但任何单一通道都不自动拥有绝对真值地位。

接口同时维护观测状态、模型状态与想象状态三类标签。观测状态来自已提交的外部数据，模型状态由预测器生成，想象状态由规划、回忆或生成模块临时构造。三者可以共享表示空间，却必须保留来源位和读写权限。内部预测不得伪造外部来源；生成候选在未通过现实反馈前不得进入“已发生事实”区。所谓意志不能改变感知、只能改变解释，也需要细化：主动视觉可以转动相机、调整曝光或触摸对象，从而改变获取何种证据；但它不能把未观测候选直接改写成观测事实。Boundary或等价边界模块负责这一来源隔离。

现实失锚可由可观测指标诊断。设标准化残差为 $r_t=W_t^{1/2}(y_t-\hat y_{t|t-1})$，传感器健康得分为 $h_t$，跨模态一致性为 $c_t^{m}$，来源覆盖率为 $q_t$。单次大残差可能代表新事件、模型错误或传感器故障，不能直接称为幻觉。我们将失锚风险定义为一个经训练或规则标定的函数
\[
R_t^{anchor}=g(\|r_t\|,h_t,c_t^{m},q_t,\text{age}_t),
\]
并要求 $g$ 的阈值在验证集上给出误报与漏报。只有当内部输出被当作外部事实提交、来源覆盖长期不足且反证未能修订状态时，才构成系统级失锚。感觉剥夺、梦境、生成模型采样和传感器断线具有不同机制，不能用狄利克雷边界转为诺伊曼边界一式概括。

内部平滑仍有合理位置。对图状态 $z$ 和固定非负关系图的拉普拉斯矩阵 $L_G$，正则项 $z^TL_Gz/2$ 鼓励相邻节点取相近表示；它帮助去噪，却不代表思维天然追求物理最小能量。当观测边界与内部预测冲突时，可求解
\[
\min_z\;\frac12\|W_o^{1/2}(Hz-y)\|^2+\frac{\lambda}{2}z^TL_Gz,
\]
其中 $H$ 是观测读出。若每个未观测连通分量都与观测或其他正定先验相连，目标严格凸时才有唯一解。若图结构错误，平滑会把错误传播得更远；因此大残差既可能要求调整 $z$，也可能要求编辑 $G$ 或重标定 $H$。模型不能用“内部连贯”压制可靠反例。

现实校准的测试应包含传感器遮挡、伪信号、时钟漂移、跨模态冲突、回放攻击与内部高置信错误。系统应能区分“没有证据”“证据支持多个解释”“证据与模型冲突”和“传感器不可用”，并采取检索、主动感知、降级或拒绝行动。现有LLM出现无依据输出，原因可能是训练目标、检索失败、解码策略、来源接口或评估缺失，不能单因没有物理传感器就归结为边界脱落。本章模型只提出一种具身现实校准接口，不把它扩展为所有幻觉现象的唯一解释。

校准结果还应按任务风险分级提交；同一置信区间足以支持低风险描述，却未必足以支持高速运动控制。阈值因任务而变，不等于事实本身因目标而改变。

\section{预测残差与复杂度门控：从“激波”到任务相关更新}
\label{sec:section_8020}

微观层必须决定哪些变化进入有限工作状态，哪些由局部闭环吸收，哪些只写入日志。历史上把预测失配称为“应力激波”、把过滤称为“麦克斯韦妖”，容易把误差、注意强度、信息熵和物理能量混为一谈。我们把上行更新定义为带类型的残差包。设观测特征为 $o_t=(F_t,Q_t,B_t)$，下行预测为 $\hat o_t=P(X_{t|t-1})$，则结构残差、属性残差和绑定残差分别为
\[
r_t^f=d_f(F_t,\hat F_t),\qquad r_t^q=d_q(Q_t,\hat Q_t),\qquad r_t^b=d_b(B_t,\hat B_t).
\]
三个度量的单位与范围分别标定，不能直接相加；若需形成总分，先以验证集尺度 $s_f,s_q,s_b$ 归一化，再用任务权重组合。结构残差高可能意味着障碍位置错误，属性残差高可能意味着颜色或材质错误，绑定残差高则可能表示“红色属于哪个对象”尚未确定。

门控器的输入不仅是惊奇，还包括任务相关性、来源可靠度、风险和处理预算。对候选更新 $e_i$，定义无量纲优先级
\[
u_i=w_r\bar r_i+w_q q_i+w_v v_i+w_d d_i-w_c c_i,
\]
其中 $\bar r_i$ 为标准化残差，$q_i$ 为当前查询相关性，$v_i$ 为安全或价值风险，$d_i$ 为截止紧迫度，$c_i$ 为预计计算成本。硬安全事件先进入资格集合，不与普通任务竞争；其余候选在容量 $C_t$ 下求解选择问题。阈值可以由TCE或其他元控制器随警觉状态调整，但调整必须记录原因和有效期。专注阅读时忽略钟摆声，是相关性和稳定背景模型共同作用；夜间行走时放大微小脚步声，是风险权重提高，而不是认知温度直接改变物理熵。

局部适应解释了衣服触压为何逐渐不进入意识。若稳定信号的均值和方差在窗口内变化很小，微观控制器可以更新背景基线 $b_t=(1-\rho)b_{t-1}+\rho y_t$，只上传偏差 $y_t-b_t$。这是一种高通或预测编码机制，但不能永久删除稳定信号：压疮、温度变化或任务要求关注衣物时，慢变量和风险检测仍需重新触发。类似地，摄像头背景减除可以降低数据量，却可能漏掉缓慢移动目标。门控策略必须用不同速度、对比度和风险场景测试，不能以“低变化等于无意义”作普遍规则。

被选中的残差包进入当前联合状态时，还要保持来源与身份。令 $J_t=\{(e_i,u_i,s_i,v_i)\}$ 为上传集合，状态更新写为
\[
X_{t+1}=F(X_t,J_t,n_t,r_t),
\]
其中 $n_t$ 是目标与硬约束，$r_t$ 是资源预算。$J_t$ 不是能量流或非齐次场源，而是经过资格检查的消息集合。它可以提高某些候选的活动、触发图编辑、请求外部工具或要求重新规划；具体效果由消息类型和更新接口决定。大残差不会自动“击碎波函数”，小残差也可能因安全关键而触发最高优先级。

CCM在AgentOS实例中可实现这一门控：输入复杂度探测器预估候选数量、耦合深度和调用成本，运行态监控器观察 $h(t)$ 的占用、冲突、延迟和读出不确定性。当负荷超限时，系统可降低非关键分辨率、把证据卸载到外部记忆 $M_e$、延迟可逆任务或请求更多资源；不可丢弃的安全信息进入保留集。一般智能系统不必使用CCM名称，但必须解决相同的有限接口问题。验证指标包括关键事件召回率、非关键抑制率、端到端延迟、负荷峰值、门控稳定性和在分布变化下的恢复能力。

门控还要避免正反馈振荡：某候选因被关注而获得更多采样，又因样本更多而持续占据注意。我们要求记录候选驻留时间、设定滞回区间，并为长期占用保留反事实抽检预算。只有当新增证据持续提高任务收益或风险判断时，优先级才可维持；否则应逐步释放容量。这样既防止新奇刺激无限霸占接口，也防止稳定但尚未解决的安全问题被一次衰减后永久遗忘。

\section{换能与采样：连续耦合、离散记录和多模态对齐}
\label{sec:section_156}

生物感受器和机器传感器都把一种物理过程转换为可供后续系统使用的信号，但“连续共振”与“离散投影”不是智能高低的本体分界。耳蜗毛细胞把机械位移转为离子通道与膜电位变化，数码相机把光子统计转成电荷并量化为像素；二者都有噪声、带宽、非线性、饱和和离散事件。我们将换能器表示为带内部状态 $s_t$ 的系统
\[
\dot s=f_s(s,x^e,t)+\xi(t),\qquad y_k=Q\!\left(\int_{t_k}^{t_k+\Delta t}g_s(s_t,x_t^e)dt\right)+\eta_k,
\]
其中 $Q$ 是量化或事件触发算子。即使输出为连续电压，后续神经脉冲或采样也引入有限分辨率；即使输出为数字，来源链、标定和高采样率也可保持重要动力学。符号接地来自可追踪的交互闭环，不要求物理场在材料中无缝延续。

采样设计必须依据任务带宽。对近似带限信号，采样率高于两倍最高相关频率可避免理想条件下的混叠；现实传感器还需抗混叠滤波、时钟同步和量化误差界。任务相关频率不是环境所有微结构的最高频率，例如抓取控制关心接触瞬态和物体运动，不需要记录每个光子的量子态。事件相机按亮度变化触发，适合高动态、低冗余视觉；帧相机保留绝对亮度和规则图像结构。神经形态传感器可以改善延迟和能效，却不会仅因脉冲形式自动获得因果理解或AGI能力。

“共振”可在工程上限定为传感器、身体和控制器在某频段内的动态耦合。设环境到传感器的传递函数为 $G_e(\omega)$，接口滤波为 $G_s(\omega)$，控制闭环为 $G_c(\omega)$，则总响应取决于乘积及反馈稳定性。相位一致可能提高定位或节律跟踪，但过强共振也可能放大噪声并导致不稳定。所谓阻抗匹配只适用于明确的物理端口和频率模型，不能推广为语义真理。系统应测量幅频响应、相位裕量、延迟和饱和边界，而不是把连续性等同于无限分辨率。

多模态对齐要求把不同采样率、坐标系和延迟统一到事件身份。相机、IMU、触觉和音频各自输出时间戳及协方差，接口通过外参变换和时间插值把候选投到共同参考系。若外参 $T_{ab}$ 变化，必须进入完整导数或作为离散重标定事件；不能假定固定几何。跨模态一致性可以降低歧义，例如视觉杯柄位置与触觉接触共同确认抓取点，但一致也可能来自共同偏差。系统应保留每个通道的独立证据，避免融合后无法追溯错误来源。

生物与机器的比较因此落在可测维度：带宽、动态范围、能耗、延迟、可塑性、故障模式和闭环适应速度。生物耳蜗并非“无限分辨率”，GPU输入也非天然与现实断裂。任何实现只要维护来源、标定、对象身份和反馈，就能成为现实锚定的一部分；任何实现若忽略这些契约，即使物理连续也可能产生错误。这一结论为HSF-HD保留实现多样性，也避免从材料差异直接推导智能本体差异。

采样接口的验证不能只用离线数据。离线回放能够比较重构误差，却不能暴露时钟漂移、缓冲拥塞、动作引起的视角变化和闭环相位延迟。我们因此同时要求静态标定、受控动态试验和真实闭环扰动：先测每个通道的噪声与延迟，再测跨模态事件能否在误差界内对齐，最后让控制动作主动改变观测并核对因果顺序。若回放性能良好而闭环不稳定，问题应归入接口时序或控制设计，不能归因于抽象表示“不够连续”。

\section{观测事件到语义状态：瞬态输入、持续表示与身份保持}
\label{sec:section_873}

历史源流 $\vec J_{ext}$ 与语义子 Semantion 的区分触及一个重要问题：瞬态输入事件与系统内部持续状态不是同一对象。我们保留这个区分，但不用“源创造粒子”的场论图景。令上传事件为
\[
e_t=(\iota_t,F_t,Q_t,B_t,S_t,\Delta_t),
\]
其中 $\iota_t$ 是对象或事件身份候选，$F_t,Q_t,B_t,S_t$ 分别是结构、属性、绑定与来源记录，$\Delta_t$ 是不确定性。事件只表示“接口在该时刻提交了什么”，持续语义状态则是经过多次证据更新后保存在 $G_t,z_t$ 和绑定表中的对象记录。一次事件可能新建节点、修订旧节点、提高或降低候选置信，也可能因无法绑定而进入待决区。

状态更新可写为类型化算子
\[
(G_{t+1},z_{t+1},B_{t+1})=U_{\tau(e_t)}(G_t,z_t,B_t,e_t),
\]
其中 $\tau(e_t)$ 是新建、属性修订、关系修订、反证、合并或拆分等事件类型。连续分布状态 $z_t$ 可以在固定结构上按微分或差分方程更新，图结构与身份变化则必须原子提交。把两者写成同一个波动方程会掩盖维度变化、来源审计和事务失败。语义状态的持续时间由任务查询、保留策略和后续证据决定，不是固定的指数衰减；重要但低活动的事实可以长期保存，高活动候选也可能在反证后立即撤销。

“打包与解包”可具体化为消息模式。包内至少包含 schema 版本、对象身份、结构字段、属性字段、来源、时间、协方差或候选概率、访问权限和完整性校验。接收端先验证模式，再把结构字段写入局部网络，把连续属性写入分布表示，把来源与版本写入证据索引。若红色与圆形来自同一感受野却存在两个邻近对象，绑定字段不能省略；若身份暂时不确定，系统保存多个假设和互斥关系。通过模式契约实现的同源绑定比狄拉克空间点更适合离散计算，也能处理遮挡、移动对象和多主体消息。

新概念形成需要累积证据与任务用途，而不是单次高强度输入越过“拓扑能隙”。设候选概念 $c$ 的支持证据集合为 $E_c$、反证为 $N_c$、查询收益为 $g_c$、维护成本为 $m_c$，结构提交器先检查最小覆盖、来源独立性和反例，再根据收益与成本决定是否建立稳定节点。强烈刺激可以提高活动和优先级，却不能代替多样证据。一次巨响可能只形成事件记录，多次在不同场景中观察可燃物与高温的关系，才可能支持可迁移规则。提交后仍保留适用域和撤销条件。

在无外部输入时，内部状态可以回忆、模拟和联想，但输出必须标为内部生成。新的外部事件到来时，观测更新与内部候选竞争；若证据不够，系统允许“不确定”而不强制坍缩为单一解释。局部结构网络提供对象身份、事件关系和来源路径，全局分布表示提供候选强度、相似性和不确定度。二者通过任务查询联合：结构相同不保证属性相同，属性相似也不保证对象同一。这个分离延续全书的形质框架，却不把其实体化为两类基本粒子。

本节的验证包括事件重放、对象跨帧保持、属性错配、合并/拆分、来源撤回与内部生成隔离。系统应能解释某个语义状态由哪些事件建立、哪些字段仍不确定、何种反例会使其撤销。这样的可审计持续表示，才是“瞬态源与稳定实体”区别的计算含义。Semantion可以继续作为书中对联合语义单元的简称，但它不是由输入能量创生的物理粒子。

\section{效应器接口：目标细化、顺应控制与安全执行}
\label{sec:section_5244}

行动不是感知编码的简单逆函数，因为从目标状态到控制序列通常是一对多问题，并受身体动力学、环境约束、授权和安全规则限制。令身体状态 $x_t^b\in\mathbb{R}^{n}$、控制输入 $u_t\in\mathbb{R}^{p}$，动力学为
\[
x_{t+1}^b=f_b(x_t^b,u_t,w_t),
\]
其中 $w_t$ 是接触与环境扰动。宏观层提供参考轨迹 $x_{t:t+L}^{ref}$、容差、硬约束集合 $\mathcal C_t$ 和任务优先序，而效应器接口在有限时域内求解合法控制。下行消息必须携带参考状态、坐标系、截止时间、最大力矩、可逆性等级和失败处理，不能只发送“向前走”或一个抽象吸引子。

对接触任务，我们可以采用阻抗或顺应控制。设位置误差 $e=x^{ref}-x^b$，在连续近似和固定模式区间内控制律为
\[
u=K_t e+D_t\dot e+u_{ff},
\]
其中 $K_t,D_t$ 分别具有刚度与阻尼的物理单位，$u_{ff}$ 是基于身体模型的前馈项。$K_t,D_t$ 必须为满足安全界的矩阵，而不是由“质元柔和/坚决”直接决定。任务语义可以通过控制策略映射器选择参数区间：探索陌生表面时降低刚度以顺应环境，搬运刚性部件时提高路径保持，但仍受力矩、碰撞和人体接触限制。语义目标先被细化为工程参数，再由底层控制器执行。

在理想固定参考、无饱和且模型满足 $M\ddot x+D\dot x+Ke=0$ 时，若 $M,K$ 正定、$D$ 正定，可选
\[
V(e,\dot e)=\frac12\dot e^TM\dot e+\frac12e^TKe,
\]
得到 $\dot V=-\dot e^TD\dot e\leq0$。这只给出该简化闭环的稳定性；时变参考、接触切换、输入饱和、估计误差和网络延迟都会增加项，需重新证明或采用鲁棒裕量。目标下降不等于抓取成功，稳定跟踪也不保证路径安全。离散实现还要声明采样周期，并验证数值积分和控制器在该步长下保持稳定。

环境阻抗估计用于调整动作，而不是追求抽象的复共轭匹配。机器人预想前方为空却碰到墙时，力传感器和视觉反馈产生结构残差，局部控制器应立即限力、停止或退让；只有无法在局部策略内消除的失配才上报宏观重规划。摸索物体、在不平地面行走或手部抖动补偿都说明微观闭环可以自主处理高频扰动。庖丁解牛的意象可理解为良好的模型、工具和顺应策略降低无效作用，并非能量传输定理。

执行资格必须先于优化。Boundary检查动作是否在授权域，安全监控器检查碰撞、力矩、速度和禁区，CCM评估控制与感知负荷；不合格动作不能因价值函数很低而执行。对于不可逆操作，系统需要双重确认、可观测反馈和中止条件。动作发出后，效应器记录实际控制、身体响应和环境结果，形成新的上行事件。若只记录指令而不记录结果，系统无法区分计划失败、执行器故障和环境改变。

因此，“逆向VTE”更准确地说是目标细化与控制接口：它从高层联合状态读取任务相关目标，输出可执行参考和参数，再以局部反馈闭环将实际轨迹保持在安全管道内。它不能保证环境被改造成内部世界图的同胚副本；现实通常只需在任务容差内达到某个状态。效应器性能由跟踪误差、响应时间、接触峰值、能耗、约束违例、恢复能力和上报质量评估。

当多个目标冲突时，接口必须显式给出优先关系，而不能把权重和为一个标量后隐藏牺牲项。安全与授权作为硬约束，任务质量和资源消耗才进入软目标；若硬约束下无可行解，控制器返回不可执行证明或近似原因，并请求上层修改目标。这个失败结果同样是有效输出，因为它阻止宏观计划把底层能力边界误读成执行意愿不足，也为后续模型修订提供了可定位证据。

\section{双向具身接口：上行证据、下行约束与闭环握手}
\label{sec:section_2301}

微观层的统一结构是双向事务接口，而不是两个纤维丛之间的全息切面。上行通道把原始信号变成带身份、来源和不确定性的观测事件；下行通道把目标和规范状态细化为参考轨迹、控制参数与资格约束；反馈通道把实际结果与预测相比较并决定局部修正还是全局重规划。三条通道共享对象身份和时间基准，任何一条缺失都会破坏闭环：无上行证据时控制变成开环，无下行约束时感知没有任务分辨率，无结果反馈时系统不能学习。

我们把接口状态写为
\[
I_t=(O_t,A_t,R_t,\Gamma_t),
\]
其中 $O_t$ 是已验证观测队列，$A_t$ 是待执行动作队列，$R_t$ 是结果和残差队列，$\Gamma_t$ 是坐标、身份、来源、权限与版本契约。上行提交 $C_o$、下行提交 $C_a$ 和反馈提交 $C_r$ 都是原子事件。宏观状态读取 $O_t$ 形成目标，效应器只有在 $C_a$ 检查通过后执行，结果再通过 $C_r$ 回写。消息超时、身份不匹配或坐标版本变化会使事务失败并触发重新估计，而不是静默沿用旧参数。

上行的任务相关性由查询决定。识别杯子时，结构网络提供杯体、杯柄、桌面和支撑关系，分布表示提供颜色、材质和可抓取性候选；抓取任务读出杯柄位姿与摩擦范围，描述任务则可能读出颜色与类别。下行不发送完整高维状态，只发送任务所需的参考与约束。这种降维不是本体丢失，而是接口契约；若任务改变，查询和下行消息也必须改变。微观层不能假定一个固定投影适用于所有目标。

闭环握手由预测、执行和确认构成。系统在动作前生成短期预测 $\hat y_{t+1:t+L}$，执行首个合法控制，接收真实 $y$ 后计算类型化残差。若残差在局部容差内，底层控制器继续；若超出容差但安全，重新估计环境和控制参数；若触及硬约束，立即中止并上报。所谓“心流”在此可描述为模型准确、负荷适中、反馈及时且局部闭环连续成功的运行区间，不能简化为物理阻抗完全匹配或零耗散。所谓“惊奇激波”则是高优先级残差消息，不是守恒能量的爆发。

局部结构网络与全局分布表示在接口处联合但不混同。前者维持“哪个传感器看见哪个对象、哪个动作作用于哪个部件、哪个结果对应哪个指令”，后者维持连续估计、候选概率和控制状态。绑定表连接两者，读出和反馈共享绑定版本。若只保留全局向量，系统容易丢失责任和来源；若只保留离散图，系统难以表示噪声、连续运动和不确定性。两者的耦合是任务相关的工程选择，不从脑区或物理场类比推出唯一机制。

AgentOS可以把 $h(t)$ 用作双向接口的当前认知总线：World、$E$、$\Theta$ 与 $M_e$ 的读出进入 $h(t)$，动作从 $h(t)$ 经Boundary和效应器发出，反馈再回到 $h(t)$。SPC提供面向主体的状态摘要，TCE调节探索和控制强度，CCM防止候选与反馈挤占有限界面，Offloading保存低频证据。一般系统也可使用其他消息总线、黑板或反应式架构；HSF-HD只要求状态、来源、约束和反馈关系可描述、可测量、可修订。

接口验收应以端到端故障注入为主：传感器延迟、错帧、对象换位、执行器饱和、工具返回乱序、目标变更和外部记忆断开。系统不仅要完成正常任务，还要在失败时保持来源、拒绝越权、停止危险动作并给出可追踪诊断。只有这种闭环完整性通过，微观层才能称为现实的看门接口。

\section{具身微观动力学总论：混合状态、稳定条件与理论位置}
\label{sec:section_5750}

我们最后把感知、状态估计、局部控制和结构事件统一成混合动力系统。连续状态记为 $x=(x^b,z)$，包含身体状态与全局分布表示；离散结构记为 $q=(G,B,\Gamma)$，包含关系图、绑定和接口模式；环境输入为 $w$，控制为 $u$，观测为 $y$。在模式 $q$ 固定的区间内，系统满足
\[
\dot x=f_q(x,u,w,t),\qquad y=h_q(x,w,t)+\nu,
\]
当验证、绑定、碰撞或模式切换条件 $g_j(x,y,q)=0$ 触发时，执行离散重置
\[
(x^+,q^+)=R_j(x^-,q^-,e_j).
\]
连续导数不能跨越 $R_j$ 直接延长；图节点增删、来源撤回和动作提交都属于离散事件。这一模型比单一哈密顿量更符合具身接口既有连续运动又有事务变化的事实。

微观控制目标由估计误差、跟踪误差、控制代价和风险组成。为避免量纲混乱，先用尺度矩阵把各项无量纲化，硬约束独立处理：
\[
J=\sum_{k=0}^{L-1}\left(\|e_k^x\|_{W_x}^2+\|e_k^y\|_{W_y}^2+\|u_k\|_{W_u}^2+\rho_k\right)+\|e_L^x\|_{W_T}^2,
\quad (x_k,u_k)\in\mathcal C_k.
\]
$\rho_k$ 是经标定的风险评分，不是焦耳能耗；实际能耗可以另以电流、电压和时间测量。优化得到的策略只在模型、时域和约束条件下有意义。环境、目标或结构改变时，$J$ 也改变，滚动控制必须重新求解。目标函数下降不能推出外部任务成功，成功由现实反馈和终止条件判定。

稳定性需要分层说明。底层固定模式控制器可用Lyapunov函数或输入到状态稳定性证明局部性质；估计器需给出可观测性、噪声和校准条件；模式切换需防止抖振和无限事件；上层重规划需保证不会长期阻塞安全反射。各层单独稳定不自动推出整体稳定，尤其当延迟、饱和和学习同时存在时。工程上应设置看门狗、最小驻留时间、速率限制、回滚模型和安全降级策略。系统在域外情形下能够停止并请求帮助，往往比继续输出一个高置信动作更重要。

学习发生在不同时间尺度。快速环更新 $h(t)$ 或等价在线状态，中速环把已验证事件写入情景记录并调整局部模型，慢速环在保留集上修改结构参数 $\Theta$ 或其他长期模型。慢速更新前后要比较现实校准、身份保持、控制稳定和旧任务性能；若任何关键指标恶化超过阈值，更新回滚。Offloading可把原始日志、工具和稀疏技能迁移到外部记忆，但安全与审计证据不能只剩不可解释参数。多尺度关系来自作用范围和更新时间，不是低频天然拥有更大权威。

本章涉及实际身体时，物理语言仍有必要：传感器带宽、质量矩阵、接触力、阻抗、焦耳能耗和耗散都是真实工程量。我们去除的是把这些量未经接口映射直接等同为语义、价值或认知证明。颜色分类概率不是能量，预测残差不是应力张量，注意门控不是麦克斯韦妖，语义单元也不是基本粒子。辅助类比可以帮助直觉，但公式的成立依赖对象、单位、假设和验证，而不依赖名称相似。

理论层级在微观层尤其清楚。系统论提供内外边界、反馈与故障隔离；协同论和复杂系统理论帮助分析多传感器、多执行器和多时间尺度耦合；耗散理论约束真实具身系统的资源供给、热与不可逆磨损；认知科学提供注意、适应、预测与行动的经验问题。在这些基础上，HSF-HD描述观测如何改变联合状态、目标如何细化为动作、反馈如何修订结构的一般状态动力学；AgentOS的 $h(t)$、$E$、$\Theta$、$\Psi$、SPC、TCE、CCM、Boundary、Offloading与 $M_e$ 是一种工程实例，而不是微观层的唯一实现。

最终验证矩阵应覆盖观测、绑定、门控、控制、学习和失效六类情形。我们要求系统在红白对象换位时不串绑，在传感器饱和时降低置信，在虎纹等突发高风险特征出现时及时上报，在衣物等稳定背景下不过载，在墙体意外接触时限力停止，在软硬物体切换时调整顺应性，在内部生成与外部观测冲突时保留来源，在模型更新后仍能完成旧任务。通过这些测试，微观层才真正承担现实锚定和行动闭环；它不是全息宇宙的物理切面，而是智能状态动力学面对现实所必须兑现的接口契约。
